\documentclass[11pt]{article}

\usepackage[T1]{fontenc}
\usepackage{lmodern}
\usepackage[a4paper,margin=28mm]{geometry}
\usepackage{amsmath,amssymb,amsthm,mathtools,mathrsfs}
\usepackage{microtype}
\usepackage{enumitem}
\usepackage{xcolor}
\usepackage[colorlinks=true,linkcolor=blue!45!black,citecolor=blue!45!black,
  urlcolor=blue!55!black]{hyperref}
\setlength{\emergencystretch}{2em}

\newtheorem{theorem}{Theorem}[section]
\newtheorem{proposition}[theorem]{Proposition}
\newtheorem{lemma}[theorem]{Lemma}
\newtheorem{corollary}[theorem]{Corollary}
\newtheorem*{maintheorem}{Main Theorem}
\theoremstyle{definition}
\newtheorem{definition}[theorem]{Definition}
\theoremstyle{remark}
\newtheorem{remark}[theorem]{Remark}

\newcommand{\C}{\mathbb C}
\newcommand{\N}{\mathbb N}
\newcommand{\cA}{\mathcal A}
\newcommand{\cD}{\mathcal D}
\newcommand{\cF}{\mathcal F}
\newcommand{\cK}{\mathcal K}
\newcommand{\cM}{\mathcal M}
\newcommand{\cP}{\mathcal P}
\newcommand{\ord}{\operatorname{ord}}
\newcommand{\supp}{\operatorname{supp}}
\newcommand{\Frac}{\operatorname{Frac}}
\newcommand{\trdeg}{\operatorname{trdeg}}

\title{Positive Arithmetic Has a M\"obius Boundary:\\
A Sign-Transmutation and Divisor-Polarity Law\\
for the Riemann Zeta Function}
\author{Michael Haufschild}
\date{}

\begin{document}

\maketitle

\begin{abstract}
We prove that signed M\"obius inversion occurs as a singular boundary of an explicit positive
arithmetic resolution.  Fix \(K\geq12\).  The \(K\)-free integers followed by successive
largest-prime-exponent shells give a complete, pairwise disjoint \(0/1\) partition
\((c_{K,j})_{j\geq0}\) of the positive integers.  The associated Dirichlet series
\(S_{K,j}=\zeta M_{K,j}\) all carry exactly the zeta divisor on
\(\Re s>\tfrac12\), with multiplicity, whereas the normalized family
\((M_{K,j})_{j\geq0}\) is differentially algebraically independent over the full differential
field generated by \(\zeta\).  Their ordered prefixes satisfy
\[
 P_{K,N}(s):=\sum_{j=0}^{N}S_{K,j}(s)
 =\frac{\zeta(s)}{\zeta((K+N)s)}.
\]
For every real \(x>1\), all native values \(P_{K,N}(x)\) are positive, but
\[
 P_{K,N}\!\left(\frac{x}{K+N}\right)
 \longrightarrow-\frac1{2\zeta(x)}<0;
\]
the critical values are eventually negative.  Their arguments have already left the common
Dirichlet half-plane of the \(0/1\) expansions.  On \(\Re u>1\), the boundary is
\[
 -\frac1{2\zeta(u)}
 =\sum_{n\geq1}\frac{-\mu(n)/2}{n^u},
\]
whose coefficients have both signs.  It therefore cannot be a coefficientwise limit of
nonnegative coefficients on a fixed Dirichlet basis.  More generally, renormalized depth
\((K+N_n)\varepsilon_n\to\lambda\) selects the boundary
\(\zeta(0)/\zeta(\lambda u)\): zero depth gives \(1\), finite nonzero depth retains a reciprocal
divisor, and right escape gives \(\zeta(0)=-1/2\).  The reciprocal boundary reverses every local
divisor order.

For the completed divisor, reflection \(s\mapsto1-s\) selects, up to scale, the unique affine
odd coordinate \(s-\tfrac12\).  Its boundary logarithmic current has residue
\(-m(\rho-\tfrac12)\) at a zero \(\rho\) of multiplicity \(m\), whereas every reflection-even
scalar observer identifies reflected partners.  Thus positive arithmetic, signed inversion,
divisor polarity, and critical-line orientation form one forced source-to-boundary chain.
As an unconditional quantitative consequence of the completed prime--Gamma source, every
off-line zero \(\rho=\beta+i\gamma\) satisfies
\(\gamma^2>7907/93\), and the total centered transverse orientation energy is
less than \(1/150000\).  No hypothesis on zeta zeros is used, and the paper does not assert RH.
\end{abstract}

\medskip
\noindent\textbf{2020 Mathematics Subject Classification.}
Primary 11M26, 11J91; Secondary 12H05, 11M41, 11A51.

\noindent\textbf{Keywords.}
Riemann zeta function; Riemann hypothesis; divisor; scale transition; hypertranscendence;
differential algebra; \(k\)-free integers; M\"obius inversion; arithmetic functions.

\section{Introduction}\label{sec:intro}

Positive Dirichlet coefficients and M\"obius inversion occupy opposite sides of a familiar
divide: the first support monotonicity, while the second supplies the signs needed for
multiplicative cancellation.  We prove that they are two scale faces of one explicit arithmetic
object.  The transition is necessarily singular: it leaves the common Dirichlet domain, changes
the sign of real values, reverses the divisor, and only then admits the affine polarization that
records signed displacement from the critical line.  This one mechanism explains why positive
coefficient arguments, unpolarized divisor counts, squared coordinates, compact fixed-scale
limits, and finite differential models can each recover part of the Riemann divisor while losing
the datum required for zero orientation.

For \(n\geq2\), write
\[
 M(n)=\max_{p\mid n}v_p(n),\qquad M(1)=0.
 \tag{1.1}
\]
Fix \(K\geq12\).  Channel zero consists of the \(K\)-free integers.  Channel \(j+1\) consists
of the integers with \(M(n)=K+j\).  If \(c_{K,j}\) denotes the corresponding indicator and
\[
 S_{K,j}(s)=\sum_{n\geq1}c_{K,j}(n)n^{-s},
 \tag{1.2}
\]
then the coefficient supports are positive, pairwise disjoint, and complete.  On
\(\Re s>1\),
\[
 S_{K,j}(s)=\zeta(s)M_{K,j}(s),
 \tag{1.3}
\]
where
\[
 M_{K,0}(s)=\frac1{\zeta(Ks)},\qquad
 M_{K,j+1}(s)=\frac1{\zeta((K+j+1)s)}
              -\frac1{\zeta((K+j)s)}.
 \tag{1.4}
\]
Every normalized multiplier is holomorphic and zero-free on
\(\Re s>\tfrac12\).  Thus all positive channels carry exactly the same zeta divisor in that
half-plane, with multiplicity.

The common divisor does not arise from hidden finite differential coupling.  Let
\[
 K_\zeta=\C(s)\langle\zeta\rangle
 =\C(s)\bigl(\zeta,\zeta',\zeta'',\ldots\bigr)
 \tag{1.5}
\]
inside a differential field of meromorphic germs.  We prove that
\[
 \{M_{K,j}^{(r)}:j,r\geq0\}
 \quad\text{is algebraically independent over }K_\zeta.
 \tag{1.6}
\]
The divisor is therefore longitudinal information shared by channels of maximal transverse
differential rank.

The new feature appears when the channels are fused in depth order.  Put
\[
 P_{K,N}(s)=\sum_{j=0}^{N}S_{K,j}(s).
 \tag{1.7}
\]
The sum is exactly
\[
 P_{K,N}(s)=\frac{\zeta(s)}{\zeta((K+N)s)}.
 \tag{1.8}
\]
The right side is the Dirichlet series of the \((K+N)\)-free integers on its half-plane of
absolute convergence, but \eqref{eq:prefix-transition} below is used as a meromorphic identity.
At fixed \(s\) in the open right half-plane, it converges to \(\zeta(s)\).  At a moving
coordinate \(s=\varepsilon_nu\to0\), its behavior is selected by the renormalized depth
\[
 \delta_n=(K+N_n)\varepsilon_n.
 \tag{1.9}
\]
If \(\delta_n\to\lambda\), then the limit is
\(\zeta(0)/\zeta(\lambda u)\), away from its poles.  If
\(\Re(\delta_nu)\to+\infty\), the limit is the constant \(\zeta(0)\).
The three vanishing-coordinate phases are consequently
\[
\begin{array}{c|c}
\text{renormalized depth}&\text{limit}\\ \hline
\delta_n\to0&1\\
\delta_n\to\lambda\ne0&\zeta(0)/\zeta(\lambda u)\\
\Re(\delta_nu)\to+\infty&\zeta(0).
\end{array}
\tag{1.10}
\]
Only the finite nonzero window retains a divisor.  In particular, exact inverse depth retains
the original divisor with opposite polarity:
\[
 P_{K,N}\!\left(\frac{u}{K+N}\right)
 \longrightarrow-\frac1{2\zeta(u)}.
 \tag{1.11}
\]
This limit changes sign, not only divisor order.  For every real \(x>1\),
\[
 P_{K,N}(x)>0\quad\hbox{for every }N,
 \qquad
 P_{K,N}\!\left(\frac{x}{K+N}\right)
 \longrightarrow-\frac1{2\zeta(x)}<0.
 \tag{1.12}
\]
The second sequence is eventually negative, but its arguments are eventually below \(1\), where
the positive Dirichlet expansions no longer converge.  On \(\Re u>1\), the limiting boundary is
\[
 -\frac1{2\zeta(u)}
 =\sum_{n\geq1}\frac{b(n)}{n^u},
 \qquad b(n)=-\frac{\mu(n)}2,
 \qquad b(1)=-\frac12<0<\frac12=b(2).
 \tag{1.13}
\]
Since pointwise limits of nonnegative coefficients remain nonnegative, no coefficientwise limit
on a fixed Dirichlet basis can produce \(b\).  The moving scale and analytic continuation are
therefore forced, not technical choices.

Completion contributes a distinct, indispensable operation.  Let
\[
 \Xi(s)=\frac12s(s-1)\pi^{-s/2}\Gamma(s/2)\zeta(s),
 \qquad \Xi(1-s)=\Xi(s).
 \tag{1.14}
\]
Among affine functions, reflection selects \(s-\tfrac12\), up to a scalar, as the unique odd
coordinate.  Hence the canonical polarized boundary current is
\[
 \mathcal J_\Xi(s)
 =-\left(s-\frac12\right)\frac{\Xi'(s)}{\Xi(s)}.
 \tag{1.15}
\]
It is reflection invariant, and at a zero \(\rho\) of multiplicity \(m\),
\[
 \operatorname*{Res}_{s=\rho}\mathcal J_\Xi
 =-m\left(\rho-\frac12\right).
 \tag{1.16}
\]
The residue contains the full signed horizontal displacement.  Squaring the affine coordinate
or taking a reflection-even norm erases that sign; integrating only logarithmic order retains
multiplicity but not displacement.  This elementary polarization becomes nondecorative here
because \eqref{eq:prefix-transition} independently produces the reciprocal zeta divisor before
completion is invoked.

\begin{maintheorem}[Positive-to-M\"obius boundary-polarity law]
\label{thm:main-intro}
For every integer \(K\geq12\), the family \((S_{K,j})_{j\geq0}\) satisfies all of the following
unconditionally.
\begin{enumerate}[label=\textup{(\alph*)}]
\item The coefficient streams are pairwise disjoint, \(0/1\)-valued, and sum coefficientwise to
the all-ones zeta stream.
\item Every channel has exactly the zeta divisor on \(\Re s>\tfrac12\), away from the pole at
\(1\), with analytic multiplicity.
\item The normalized channels \((M_{K,j})_{j\geq0}\) are jointly differentially algebraically
independent over \(K_\zeta\).
\item The exact finite-depth law \eqref{eq:prefix-transition} holds, and its fixed, finite
renormalized-depth, and right-escaping limits are precisely those in
\eqref{eq:scale-selector-theorem}--\eqref{eq:canonical-scale-witnesses}.
\item For every real \(x>1\), \(P_{K,N}(x)>0\) for all \(N\), whereas the critically rescaled
values in \eqref{eq:positive-mobius-sign} are eventually negative.  Their limit is the
mixed-sign M\"obius series \eqref{eq:mobius-boundary-series}; it cannot be a coefficientwise
limit of nonnegative coefficients on a fixed Dirichlet basis.
\item The exceptional boundary reverses local divisor order.  Reflection completion has the
unique affine polarization \eqref{eq:polarized-current},
whose residue is the full displacement vector \eqref{eq:polarized-residue}.
\item If \(\rho=\beta+i\gamma\) is an off-line nontrivial zero, then
\[
 \gamma^2>\frac{7907}{93}.
 \tag{1.17}
\]
The completed transverse orientation energy \(\mathcal E_\perp\), defined intrinsically in
Section~\ref{sec:orientation-budget}, satisfies
\[
 0\leq\mathcal E_\perp
 <\frac1{150000}.
 \tag{1.18}
\]
\end{enumerate}
\end{maintheorem}

The algebraic mechanism behind part~\textup{(c)} is an
\(\Omega\)-leading-echelon law.  Prime derivations lower a shell's least total-prime-factor
grade to a pure prime ray; finite-jet Jacobians become products of prime-log Vandermonde
determinants, and higher grades cannot cancel them.  Parts~\textup{(d)}--\textup{(f)}
identify what that transverse theorem cannot see: a singular source-scale limit creates signed
inversion, the exceptional boundary reverses divisor polarity, and reflection polarization
carries orientation.

The individual ingredients have classical ancestors.  The \(k\)-free identity
\(\sum\mathbf1_{k\text{-free}}(n)n^{-s}=\zeta(s)/\zeta(ks)\) is standard and remains active in
work on \(k\)-free error terms and analytic continuation
\cite{Meng2015,Kawalec2026}.  Hypertranscendence of Dirichlet series goes back to
Ostrowski~\cite{Ostrowski1920}; differential and algebraic independence were developed by
Postnikov, Voronin, Shapiro--Sparer, and Pong
\cite{Postnikov1949,Voronin1973,ShapiroSparer1986,Pong2016}.  Zero-preserving zeta modifications were
constructed by Grosswald and Schnitzer~\cite{GrosswaldSchnitzer1978}.  These results do not
supply the necessity law proved here: one complete positive disjoint resolution of maximal
relative differential rank whose finite prefixes change sign at precisely the moving scale that
exposes mixed-sign M\"obius inversion, reverse every divisor order, and admit a unique affine
reflection polarization.

This is progress toward RH, not a proof of RH.  It closes the structural diagnosis: positive
bulk data can reach the inverse divisor only through singular scale escape; inverse divisor data
can retain horizontal side only through odd affine polarization; every even scalarization has
already discarded that side.  Part~\textup{(g)} and \eqref{eq:orientation-cubic} quantify the
remaining transverse defect without claiming its vanishing.

Sections~\ref{sec:shells}--\ref{sec:application} prove positive fusion, the zero-free
multiplier law, and maximal relative differential rank.
Section~\ref{sec:rhframe} forms the complete channel resolution.
Section~\ref{sec:scale-transition} proves the scale-selector and forced sign-transmutation laws.
Section~\ref{sec:polarization} proves the universal affine polarization and classifies the
principal information-losing projections.
Section~\ref{sec:orientation-budget} gives the unconditional completed-source bounds and the
exact scope of the RH consequence.

\section{Exponent shells and exact positive fusion}\label{sec:shells}

For \(r\geq1\), call \(n\) \(r\)-free if no \(p^r\) divides \(n\), and write
\[
 b_r(n)=\mathbf 1_{\{n\ {\rm is}\ r\text{-free}\}}.
\]
The convention makes \(b_1=\mathbf 1_{\{1\}}\).  Since \(r\)-freeness implies
\((r+1)\)-freeness,
\[
 a_m=b_{m+1}-b_m
 \label{eq:shell-difference}
 \tag{2.1}
\]
takes only the values zero and one.

\begin{proposition}[Positive shell partition]\label{prop:partition}
For every \(n\geq1\) and \(m,k\geq1\):
\begin{enumerate}[label=\textup{(\roman*)}]
\item \(a_m(n)=1\) if and only if \(M(n)=m\);
\item if \(m\ne k\), then \(a_m(n)a_k(n)=0\);
\item
\[
 \sum_{m\geq1}a_m(n)=1-\mathbf 1_{\{1\}}(n),
 \label{eq:pointwise-fusion}
 \tag{2.2}
\]
and the sum has finite support for each fixed \(n\);
\item for every \(R\geq1\),
\[
 \sum_{m=1}^{R}a_m(n)=b_{R+1}(n)-b_1(n).
 \label{eq:finite-telescope}
 \tag{2.3}
\]
\end{enumerate}
\end{proposition}

\begin{proof}
If \(n\geq2\), then \(n\) is \((m+1)\)-free but not \(m\)-free exactly when the largest
valuation of \(n\) is \(m\).  This proves (i) and (ii).  Formula \eqref{eq:finite-telescope}
is the telescoping sum of \eqref{eq:shell-difference}.  Taking
\(R\geq M(n)\) proves \eqref{eq:pointwise-fusion}; for \(n=1\), every difference is zero.
\end{proof}

The Dirichlet series of \(b_r\) is obtained directly from Euler factors:
\[
 \sum_{n\geq1}\frac{b_r(n)}{n^s}
 =\prod_p(1+p^{-s}+\cdots+p^{-(r-1)s})
 =\frac{\zeta(s)}{\zeta(rs)}
 \qquad(\Re s>1),
 \label{eq:mfree-series}
 \tag{2.4}
\]
where for \(r=1\) both sides equal one.

\begin{proposition}[Shell factorization]\label{prop:factor}
For every \(m\geq1\) and \(\Re s>1\),
\begin{equation}\label{eq:factor}
 S_m(s)=\zeta(s)H_m(s),\qquad
 H_m(s)=\frac1{\zeta((m+1)s)}-\frac1{\zeta(ms)}.
 \tag{2.5}
\end{equation}
Moreover the coefficient fusion in Proposition~\ref{prop:partition} is absolutely summable on
every half-plane \(\Re s\geq1+\varepsilon\).
\end{proposition}

\begin{proof}
Subtract \eqref{eq:mfree-series} at \(r=m+1\) and \(r=m\).  For
\(\sigma=\Re s>1\),
\[
 \sum_{m\geq1}\sum_{n\geq1}|a_m(n)n^{-s}|
 =\sum_{n\geq2}n^{-\sigma}<\infty,
\]
which proves the final assertion.
\end{proof}

In particular the countably infinite relation
\[
 \sum_{m\geq1}S_m(s)=\zeta(s)-1,\qquad \Re s>1,
 \label{eq:analytic-fusion}
 \tag{2.6}
\]
is exact.  The differential-independence theorem below concerns finite differential
polynomials, so \eqref{eq:analytic-fusion} is compatible with maximal finite algebraic rank:
an infinite locally convergent sum is not a differential-polynomial relation.

\section{A zero-free adjacent-dilation law}\label{sec:zerofree}

The factor \(H_m\) is analytic on \(\Re s>1/2\) for \(m\geq12\), because both dilated zeta
arguments then lie in \(\Re z>6\).  Its nonvanishing is equivalent to
\[
 \zeta(ms)-\zeta((m+1)s)\ne0.
\]
The following quantitative estimate proves more than is needed.

\begin{lemma}[Prime-two domination]\label{lem:two-dominates}
Let \(m\geq12\), \(s=\sigma+it\), and \(\sigma>1/2\).  Put
\[
 A=2^{-ms}-2^{-(m+1)s},
\qquad
 T=\sum_{n\geq3}\bigl(n^{-ms}-n^{-(m+1)s}\bigr).
\]
Then
\[
 |A|>2^{-m\sigma}\frac{29}{100},
 \qquad
 |T|<2^{-m\sigma}\frac{21}{100}.
 \tag{3.1}
\]
\end{lemma}

\begin{proof}
The reverse triangle inequality and \(|2^{-s}|=2^{-\sigma}<2^{-1/2}<71/100\) give
\[
 |A|=2^{-m\sigma}|1-2^{-s}|
 >2^{-m\sigma}\frac{29}{100}.
\]
Set \(x=m\sigma>6\).  For \(n\geq3\),
\[
 |n^{-ms}-n^{-(m+1)s}|
 \leq n^{-x}(1+n^{-\sigma})
 <\frac85 n^{-x}.
\]
Since \(n\geq3\) and \(x>6\),
\[
 n^{-x}=2^{-x}(2/n)^x
 \leq2^{-x}\frac{64}{n^6}.
\]
Consequently
\[
 |T|
 \leq2^{-m\sigma}\frac{512}{5}\sum_{n\geq3}\frac1{n^6}.
\]
The elementary integral estimate
\[
 \sum_{n\geq3}\frac1{n^6}
 \leq\frac1{3^6}+\frac1{4^6}+\frac1{2500}
\]
gives
\[
 \frac{512}{5}
 \left(\frac1{729}+\frac1{4096}+\frac1{2500}\right)
 <\frac{21}{100},
\]
proving the claim.
\end{proof}

\begin{theorem}[Zero-free shell multiplier]\label{thm:zero-free}
For every integer \(m\geq12\), \(H_m\) is holomorphic and zero-free on
\[
 \cD_{1/2}=\{s\in\C:\Re s>1/2\}.
\]
\end{theorem}

\begin{proof}
Absolute convergence gives
\[
 \zeta(ms)-\zeta((m+1)s)=A+T.
\]
Lemma~\ref{lem:two-dominates} implies \(|A|>|T|\), so this difference is nonzero.
Furthermore \(\zeta(ms)\) and \(\zeta((m+1)s)\) are holomorphic and nonzero because their
arguments have real part greater than one.  The identity
\[
 H_m(s)=
 \frac{\zeta(ms)-\zeta((m+1)s)}
 {\zeta(ms)\zeta((m+1)s)}
\]
finishes the proof.
\end{proof}

\begin{corollary}[Exact divisor replication]\label{cor:divisor}
Let \(m\geq12\).  The meromorphic continuation
\[
 S_m^\#(s)=\zeta(s)H_m(s)
\]
has, at every \(s\in\cD_{1/2}\setminus\{1\}\), exactly the same vanishing order as \(\zeta\):
\[
 \ord_s S_m^\#=\ord_s\zeta.
 \tag{3.2}
\]
In particular their zero sets in \(\cD_{1/2}\setminus\{1\}\) agree, with analytic
multiplicities.
\end{corollary}

\begin{proof}
The order of a product of holomorphic functions is the sum of their orders.
Theorem~\ref{thm:zero-free} gives \(\ord_sH_m=0\).
\end{proof}

The integer \(12\) is a uniform threshold delivered by the explicit estimate; no optimality is
claimed.  Its role is to place \(m\Re s\) strictly beyond six throughout \(\cD_{1/2}\).

\section{The \texorpdfstring{\(\Omega\)}{Omega}-leading-echelon law}\label{sec:omega}

We now isolate the general algebraic mechanism.  Let \(\cA\) be the \(\C\)-algebra of arithmetic
functions on the positive integers with Dirichlet convolution
\[
 (f*g)(n)=\sum_{d\mid n}f(d)g(n/d).
\]
It is an integral domain: if \(r\) and \(t\) are the least indices at which nonzero \(f\) and
\(g\) occur, then \((f*g)(rt)=f(r)g(t)\ne0\).

For a prime \(q\), define the prime derivation
\[
 (\partial_qf)(n)=v_q(nq)f(nq).
 \label{eq:prime-derivation}
 \tag{4.1}
\]
A coefficient check gives
\[
 \partial_q(f*g)=(\partial_qf)*g+f*(\partial_qg).
 \tag{4.2}
\]
For \(k\geq0\), write
\[
 L^kf(n)=(-\log n)^k f(n).
 \tag{4.3}
\]
Under Dirichlet transformation, \(L^k\) becomes the \(k\)-th ordinary derivative.

\begin{theorem}[\(\Omega\)-leading-echelon criterion]\label{thm:omega}
Let \(R\) be a finite nonempty set, let \(\ell:R\to\N_{>0}\) be injective, and let
\(f_r\in\cA\).  Suppose that for every \(r\in R\), prime \(q\), and \(k,n\geq0\),
\begin{align}
 \Omega(n)<\ell(r)-1
 &\ \Longrightarrow\ 
 \partial_q(L^kf_r)(n)=0,
 \label{eq:below}
 \tag{4.4}\\
 \Omega(n)=\ell(r)-1
 &\ \Longrightarrow\
 \partial_q(L^kf_r)(n)
 =\ell(r)(-\ell(r)\log q)^k
 \mathbf 1_{\{q^{\ell(r)-1}\}}(n).
 \label{eq:leading}
 \tag{4.5}
\end{align}
Then for every \(J\geq0\), the rectangle
\[
 \{L^kf_r:r\in R,\ 0\leq k\leq J\}
 \label{eq:jet-rectangle}
 \tag{4.6}
\]
is algebraically independent over \(\C\) in \(\cA\).

If, in addition, all \(f_r\) have a finite half-plane of absolute Dirichlet convergence, then
their Dirichlet transforms \(F_r\) are jointly differentially algebraically independent over
\(\C(s)\): the set \(\{F_r^{(k)}:r\in R,k\geq0\}\) is algebraically independent over
\(\C(s)\).
\end{theorem}

\begin{proof}
Fix \(J\), and index the columns by
\[
 I=R\times\{0,\ldots,J\}.
\]
For each \(i=(r,t)\in I\), choose a distinct prime \(q_i\), and form the
\(\cA\)-valued Jacobian
\[
 \mathcal J_{i,(u,k)}
 =\partial_{q_i}(L^kf_u).
\]
Put \(e(u,k)=\ell(u)-1\) and
\[
 N=\prod_{i=(r,t)\in I}q_i^{e(r,t)}.
\]
Every convolution term contributing to \(\det(\mathcal J)(N)\) has total \(\Omega\)-grade
\(\sum_i e(i)\).  By \eqref{eq:below}, no factor can occur below its column grade.  Hence every
factor must lie in its lowest possible grade, and \eqref{eq:leading} computes the coefficient
exactly.

Unique factorization forces a contributing permutation to preserve the equal-\(e\) blocks.
Because \(\ell\) is injective, these are precisely the blocks indexed by one \(u\in R\).
Up to a nonzero scalar, the \(u\)-block is
\[
 \left[(-\ell(u)\log q_{(u,t)})^k\right]_
 {0\leq t,k\leq J}.
\]
Its determinant is a Vandermonde product
\[
 (-\ell(u))^{J(J+1)/2}
 \prod_{0\leq t<v\leq J}
 \bigl(\log q_{(u,v)}-\log q_{(u,t)}\bigr),
\]
which is nonzero because the chosen primes are distinct.  Therefore
\(\det(\mathcal J)(N)\ne0\), so \(\det(\mathcal J)\ne0\) in \(\cA\).

The Jacobian criterion over a characteristic-zero domain now implies algebraic independence:
if a nonzero polynomial relation of least total degree existed, applying the derivations would
place all its nonzero partial derivatives in the kernel, contradicting invertibility of the
Jacobian over \(\Frac(\cA)\).  This proves \eqref{eq:jet-rectangle}.

It remains to justify the coordinate field \(\C(s)\).  The Dirichlet transform is an injective
algebra homomorphism from \((\cA,*)\) into meromorphic germs whenever the sources have finite
abscissa, and it sends \(L^kf\) to the \(k\)-th derivative of the transform.  Suppose, after
clearing denominators, that
\[
 P\bigl(s;(F_r^{(k)}(s))_{r,k}\bigr)=0,
 \qquad P\in\C[s][X_{r,k}]
 \label{eq:coordinate-relation}
 \tag{4.7}
\]
is a nonzero relation involving a finite jet rectangle.  Expand
\(P=\sum_{d=0}^Ds^dP_d\).  Dirichlet multiplication rewrites \eqref{eq:coordinate-relation} as
\[
 \sum_{n\geq1}Q_n(s)n^{-s}=0,
 \qquad
 Q_n(s)=\sum_{d=0}^Ds^d
 \bigl(P_d((L^kf_r)_{r,k})\bigr)(n).
 \tag{4.8}
\]
Uniqueness remains valid for Dirichlet series with uniformly bounded-degree polynomial
coefficients: choose the least \(n\) for which \(Q_n\ne0\), multiply by \(n^s\), and let
\(s=\sigma\to+\infty\); exponential separation dominates the fixed polynomial growth and forces
\(Q_n=0\), a contradiction.  Thus every \(Q_n\) vanishes identically.  Comparing powers of \(s\)
shows \(P_d((L^kf_r)_{r,k})=0\) in \(\cA\) for every \(d\).  The first part forces every
\(P_d=0\), contrary to the choice of \(P\).  Hence no relation over \(\C(s)\) exists.
\end{proof}

\begin{remark}
The proof uses the lowest nonzero associated-graded layer only.  Arbitrary higher
\(\Omega\)-support may be added to each source without changing the conclusion.  This
``upper-tail immunity'' is what makes the criterion applicable to full exponent shells rather
than only to isolated prime powers.
\end{remark}

\section{Zeta and the shells have maximal relative rank}\label{sec:application}

We verify the hypotheses of Theorem~\ref{thm:omega}.  The two elementary facts below are the
entire arithmetic input.

\begin{lemma}[Lowest shell grade]\label{lem:lowest-shell}
If \(a_m(n)\ne0\), then
\[
 \Omega(n)\geq m.
\]
If equality holds, then \(n=p^m\) for a unique prime \(p\).
\end{lemma}

\begin{proof}
Because \(a_m(n)\ne0\), some \(p^m\) divides \(n\), giving \(\Omega(n)\geq m\).
Write \(n=p^mu\).  Additivity of \(\Omega\) gives
\(\Omega(n)=m+\Omega(u)\).  Equality forces \(\Omega(u)=0\), hence \(u=1\).
Uniqueness follows from unique factorization.
\end{proof}

\begin{lemma}[Prime-ray leading layer]\label{lem:prime-ray}
For every \(m\geq1\), prime \(q\), and \(k,n\geq0\),
\begin{align*}
 \Omega(n)<m-1
 &\ \Longrightarrow\
 \partial_q(L^ka_m)(n)=0,\\
 \Omega(n)=m-1
 &\ \Longrightarrow\
 \partial_q(L^ka_m)(n)
 =m(-m\log q)^k\mathbf 1_{\{q^{m-1}\}}(n).
\end{align*}
\end{lemma}

\begin{proof}
If the left side is nonzero, then \(a_m(nq)\ne0\), and
Lemma~\ref{lem:lowest-shell} gives
\[
 m\leq\Omega(nq)=\Omega(n)+1.
\]
This proves the first assertion.  Under equality \(\Omega(n)=m-1\), nonvanishing forces
\(\Omega(nq)=m\), so Lemma~\ref{lem:lowest-shell} gives \(nq=p^m\).  Since \(q\mid p^m\),
one has \(p=q\) and \(n=q^{m-1}\).  At that point
\[
 v_q(q^m)(-\log q^m)^ka_m(q^m)
 =m(-m\log q)^k.
\]
\end{proof}

Let \(\mathbf 1(n)=1\) be the zeta coefficient source.  Assign it level one.  If
\(\Omega(n)=0\), then \(n=1\), and
\[
 \partial_q(L^k\mathbf1)(n)
 =(-\log q)^k\mathbf1_{\{1\}}(n),
\]
which is exactly \eqref{eq:leading} with \(\ell=1\).  There is no lower grade.

\begin{theorem}[Zeta--shell differential freeness]\label{thm:zeta-shell-free}
The countable family
\[
 \{\zeta\}\cup\{S_m:m\geq2\}
\]
is jointly differentially algebraically independent over \(\C(s)\).  Equivalently,
\[
 \{\zeta^{(k)}:k\geq0\}
 \cup
 \{S_m^{(k)}:m\geq2,\ k\geq0\}
 \tag{5.1}
\]
is algebraically independent over \(\C(s)\).
\end{theorem}

\begin{proof}
Take an arbitrary finite subfamily.  Assign level \(1\) to the zeta source \(\mathbf1\) and
level \(m\) to \(a_m\).  The levels are positive and distinct.  The preceding calculation
handles level one, while Lemma~\ref{lem:prime-ray} handles every shell.  All sources are bounded,
so their Dirichlet series converge absolutely on \(\Re s>1\).  Theorem~\ref{thm:omega} applies.
Because every alleged differential-polynomial relation uses only finitely many functions and
derivatives, the finite conclusion proves the countable statement.
\end{proof}

We now separate the common divisor from its transverse directions.  We use the following
standard differential-field observation.

\begin{lemma}[Base adjunction and gauge invariance]\label{lem:gauge}
Let \(K\) be a differential field of characteristic zero.
\begin{enumerate}[label=\textup{(\roman*)}]
\item If \(\{z\}\cup\{y_i:i\in I\}\) is differentially algebraically independent over \(K\),
then \(\{y_i:i\in I\}\) is differentially algebraically independent over
\(K\langle z\rangle\).
\item If \(\{y_i:i\in I\}\) is differentially algebraically independent over \(K\) and
\(a\in K^\times\), then \(\{ay_i:i\in I\}\) is differentially algebraically independent over
\(K\).
\end{enumerate}
\end{lemma}

\begin{proof}
For (i), the set of all derivatives of the \(y_i\) is algebraically independent over the
rational-function field generated by all derivatives of \(z\), by the ordinary tower property
of algebraic independence.

For (ii), for each \(k\),
\[
 (ay_i)^{(k)}
 =a\,y_i^{(k)}
 +\sum_{j<k}\binom{k}{j}a^{(k-j)}y_i^{(j)}.
\]
This is an invertible triangular change of jet coordinates over \(K\), with diagonal entry
\(a\ne0\).
\end{proof}

\begin{theorem}[Relative multiplier hypertranscendence]\label{thm:relative}
Let
\[
 K_\zeta=\C(s)\langle\zeta\rangle.
\]
Then the family
\[
 \{H_m:m\geq2\}
\]
is jointly differentially algebraically independent over \(K_\zeta\).  Equivalently, every
finite subset of
\[
 \{H_m^{(k)}:m\geq2,\ k\geq0\}
 \tag{5.2}
\]
is algebraically independent over \(K_\zeta\).
\end{theorem}

\begin{proof}
Theorem~\ref{thm:zeta-shell-free} and Lemma~\ref{lem:gauge}(i) show that the \(S_m\) are jointly
differentially algebraically independent over \(K_\zeta\).  Since \(\zeta^{-1}\) is a nonzero
element of this meromorphic differential field, Lemma~\ref{lem:gauge}(ii) shows the same for
\(\zeta^{-1}S_m\).  On the initial half-plane of convergence,
\(\zeta^{-1}S_m=H_m\) by Proposition~\ref{prop:factor}; hence the meromorphic germs agree.
\end{proof}

\begin{corollary}[Exact differential rank]\label{cor:rank}
For distinct \(m_1,\ldots,m_r\geq2\) and integers \(J_1,\ldots,J_r\geq0\),
\[
 \trdeg_{K_\zeta}
 K_\zeta\bigl(H_{m_i}^{(k)}:
 1\leq i\leq r,\ 0\leq k\leq J_i\bigr)
 =\sum_{i=1}^r(J_i+1).
 \tag{5.3}
\]
\end{corollary}

\begin{proof}
The displayed generators are algebraically independent by Theorem~\ref{thm:relative}.
\end{proof}

Corollary~\ref{cor:rank} makes ``maximal'' literal: every new shell and every new derivative
raises transcendence degree by exactly one, even after the complete differential data of zeta
have been placed in the coefficient field.

\section{The complete positive divisor frame}\label{sec:rhframe}

The preceding shell family omits the integer \(1\) and all levels below a chosen zero-free
threshold.  We now add one positive head and obtain a complete resolution.

For \(K\geq2\), let \(b_K\) be the \(K\)-free indicator from
Section~\ref{sec:shells}, and define
\[
 c_{K,0}=b_K,\qquad c_{K,j+1}=a_{K+j}\quad(j\geq0).
 \tag{6.1}
\]
Set
\[
 M_{K,0}(s)=\zeta(Ks)^{-1},\qquad
 M_{K,j+1}(s)=H_{K+j}(s).
 \label{eq:complete-multipliers}
 \tag{6.2}
\]

\begin{proposition}[Complete positive resolution]\label{prop:complete-resolution}
For every \(K\geq2\):
\begin{enumerate}[label=\textup{(\roman*)}]
\item \(c_{K,j}(n)\in\{0,1\}\);
\item \(c_{K,i}(n)c_{K,j}(n)=0\) for \(i\ne j\);
\item for every positive integer \(n\),
\[
 \sum_{j\geq0}c_{K,j}(n)=1,
 \tag{6.3}
\]
and the sum is locally finite;
\item on \(\Re s>1\),
\[
 S_{K,j}(s):=\sum_{n\geq1}c_{K,j}(n)n^{-s}
 =\zeta(s)M_{K,j}(s).
 \tag{6.4}
\]
\end{enumerate}
\end{proposition}

\begin{proof}
If \(M(n)<K\), including \(n=1\), then \(n\) lies in the \(K\)-free head and in no shell
\(a_{K+j}\).  If \(M(n)\geq K\), there is a unique \(j=M(n)-K\) for which
\(a_{K+j}(n)=1\), and \(n\) is not \(K\)-free.  This proves the first three statements.
Formula \eqref{eq:mfree-series} gives the head factorization; Proposition~\ref{prop:factor}
gives every remaining channel.
\end{proof}

The differential statement for the head is not a consequence of
Theorem~\ref{thm:relative} by itself.  We record the stronger tail law.

\begin{proposition}[Reciprocal-dilation tail]\label{prop:reciprocal-tail}
For every \(K\geq2\), the family
\[
 Y_{K+j}(s)=\zeta((K+j)s)^{-1},\qquad j\geq0,
 \tag{6.5}
\]
is jointly differentially algebraically independent over \(K_\zeta\).
\end{proposition}

\begin{proof}
Let \(\mu_m\) be the arithmetic function supported on perfect \(m\)-th powers, with
\(\mu_m(r^m)=\mu(r)\).  Its Dirichlet transform is \(Y_m\).  If
\(\partial_q\) is the prime derivation from \eqref{eq:prime-derivation}, then below
\(\Omega\)-grade \(m-1\), \(\partial_qL^k\mu_m\) vanishes.  In grade \(m-1\), its only
nonzero coefficient occurs at \(q^{m-1}\) and equals
\[
 -m(-m\log q)^k.
 \tag{6.6}
\]
Thus the \(\Omega\)-leading-echelon proof of Theorem~\ref{thm:omega} applies simultaneously to
the zeta coefficient source at level \(1\) and any finite collection of the
\(\mu_m\), \(m\geq2\).  It follows that
\(\{\zeta\}\cup\{Y_m:m\geq K\}\) is jointly differentially algebraically independent over
\(\C(s)\).  Base adjunction, Lemma~\ref{lem:gauge}(i), gives the claimed independence over
\(K_\zeta\).
\end{proof}

\begin{theorem}[Complete positive hypertranscendental divisor frame]
\label{thm:complete-frame}
Let \(K\geq12\).  Then:
\begin{enumerate}[label=\textup{(\roman*)}]
\item \((c_{K,j})_{j\geq0}\) is a complete pairwise-disjoint \(0/1\) coefficient resolution;
\item every \(M_{K,j}\) is holomorphic and zero-free on \(\Re s>\tfrac12\);
\item every continued channel \(\zeta M_{K,j}\) has exactly the zeta divisor there, away from
\(s=1\), with multiplicity;
\item \((M_{K,j})_{j\geq0}\) is jointly differentially algebraically independent over
\(K_\zeta\).
\end{enumerate}
\end{theorem}

\begin{proof}
Part~\textup{(i)} is Proposition~\ref{prop:complete-resolution}.  For the head,
\(\Re(Ks)>6\), so \(\zeta(Ks)^{-1}\) is holomorphic and nonzero.  For the shells, use
Theorem~\ref{thm:zero-free}.  The product rule for analytic order proves
part~\textup{(iii)}.

By Proposition~\ref{prop:reciprocal-tail}, the \(Y_{K+j}\) are jointly differentially
algebraically independent over \(K_\zeta\).  The transformation
\[
 Y_K,\ Y_{K+1},Y_{K+2},\ldots
 \longmapsto
 Y_K,\ Y_{K+1}-Y_K,\ Y_{K+2}-Y_{K+1},\ldots
 \tag{6.7}
\]
is invertible and triangular on every finite jet packet.  It therefore preserves algebraic
independence and gives part~\textup{(iv)}.
\end{proof}

This theorem is the first half of the main law.  Transverse complexity and longitudinal divisor
complexity decouple maximally: no integer coefficient is shared, no finite differential
relation is shared after normalization, yet the entire divisor is shared.

\section{The divisor scale-transition law}\label{sec:scale-transition}

The ordered resolution contains additional information that an unordered frame does not.  By
telescoping \eqref{eq:complete-multipliers},
\begin{equation}
 P_{K,N}(s):=\sum_{j=0}^{N}S_{K,j}(s)
 =\zeta(s)\sum_{j=0}^{N}M_{K,j}(s)
 =\frac{\zeta(s)}{\zeta((K+N)s)}.
 \label{eq:prefix-transition}
\end{equation}
This holds as a meromorphic identity.  Write more generally
\[
 A_m(s)=\frac{\zeta(s)}{\zeta(ms)}.
 \tag{7.2}
\]

\begin{theorem}[Scale selector]\label{thm:scale-selector}
Let \(m_n\) be positive integers, let \(\varepsilon_n\in\C\), and suppose
\(\varepsilon_n\to0\).  Fix \(u,\lambda\in\C\).
\begin{enumerate}[label=\textup{(\roman*)}]
\item If \(m_n\varepsilon_n\to\lambda\), \(\lambda u\ne1\), and
\(\zeta(\lambda u)\ne0\), then
\begin{equation}
 A_{m_n}(\varepsilon_nu)
 \longrightarrow
 \frac{\zeta(0)}{\zeta(\lambda u)}
 =-\frac1{2\zeta(\lambda u)}.
 \label{eq:scale-selector-theorem}
\end{equation}
\item If
\(\Re(m_n\varepsilon_nu)\to+\infty\), then
\[
 A_{m_n}(\varepsilon_nu)\longrightarrow\zeta(0)=-\frac12.
 \tag{7.4}
\]
\item At fixed \(s\) with \(\Re s>0\),
\[
 A_m(s)\longrightarrow\zeta(s)\qquad(m\to\infty).
 \tag{7.5}
\]
\end{enumerate}
\end{theorem}

\begin{proof}
For part~\textup{(i)},
\[
 A_{m_n}(\varepsilon_nu)
 =\frac{\zeta(\varepsilon_nu)}
        {\zeta((m_n\varepsilon_n)u)}.
 \tag{7.6}
\]
The numerator tends to \(\zeta(0)=-\tfrac12\).  The denominator tends to
\(\zeta(\lambda u)\) by continuity away from the pole at \(1\), and its limiting value is
nonzero, so division proves the claim.

For part~\textup{(ii)}, the numerator has the same limit.  If
\(\Re w>1\), the absolutely convergent Dirichlet series gives
\[
 |\zeta(w)-1|
 \leq\sum_{r\geq2}r^{-\Re w}
 \leq\frac1{\Re w-1}.
 \tag{7.7}
\]
Taking \(w=(m_n\varepsilon_n)u\) proves that the denominator tends to one.
For part~\textup{(iii)}, apply the same estimate to \(w=ms\).
\end{proof}

\begin{corollary}[Three vanishing-coordinate phases]\label{cor:three-phases}
Away from poles of the displayed meromorphic limits,
\begin{align}
 A_{n+1}\!\left(\frac{u}{n+1}\right)
 &\longrightarrow-\frac1{2\zeta(u)},\notag\\
 A_{(n+1)^2}\!\left(\frac{u}{n+1}\right)
 &\longrightarrow-\frac12\qquad(\Re u>0),\notag\\
 A_{n+1}\!\left(\frac{u}{(n+1)^2}\right)
 &\longrightarrow1.
 \label{eq:canonical-scale-witnesses}
\end{align}
Every sufficiently large \(A_m\) is a finite prefix \(P_{K,m-K}\).
\end{corollary}

\begin{proof}
The first and third limits are Theorem~\ref{thm:scale-selector} with
\(\lambda=1\) and \(\lambda=0\), respectively.  In the second,
\(m_n\varepsilon_n=n+1\), so the renormalized argument escapes through the right half-plane.
The prefix assertion is \eqref{eq:prefix-transition}.
\end{proof}

\subsection{Forced positive-to-M\"obius sign transmutation}

\begin{theorem}[Positive-to-M\"obius transmutation]\label{thm:positive-mobius}
Fix \(K\geq12\).
\begin{enumerate}[label=\textup{(\roman*)}]
\item For every real \(x>1\) and every \(N\geq0\),
\[
 P_{K,N}(x)>0.
 \]
\item For every real \(x>1\),
\begin{equation}
 P_{K,N}\!\left(\frac{x}{K+N}\right)
 \longrightarrow-\frac1{2\zeta(x)}<0,
 \qquad
 P_{K,N}\!\left(\frac{x}{K+N}\right)<0
 \quad\hbox{eventually}.
 \label{eq:positive-mobius-sign}
 \tag{7.9}
\end{equation}
\item The critical arguments eventually leave the common positive Dirichlet domain:
\[
 \frac{x}{K+N}<1\qquad\hbox{eventually}.
 \tag{7.10}
\]
\item On \(\Re u>1\),
\begin{equation}
 -\frac1{2\zeta(u)}
 =\sum_{n\geq1}\frac{b(n)}{n^u},
 \qquad
 b(n)=-\frac{\mu(n)}2,
 \qquad
 b(1)=-\frac12<0<\frac12=b(2).
 \label{eq:mobius-boundary-series}
 \tag{7.11}
\end{equation}
\item There is no family \(a_r(n)\geq0\) on a fixed Dirichlet basis such that
\[
 \lim_{r\to\infty}a_r(n)=b(n)
 \qquad\hbox{for every }n.
 \]
\end{enumerate}
\end{theorem}

\begin{proof}
Part~\textup{(i)} follows either from the \((K+N)\)-free Dirichlet series, whose coefficients
are nonnegative and whose coefficient at \(1\) is one, or from positivity of
\(\zeta(x)/\zeta((K+N)x)\) for \(x>1\).

Part~\textup{(ii)} is the first limit in
\eqref{eq:canonical-scale-witnesses}.  The real value \(\zeta(x)\) is positive for \(x>1\);
hence the limit is strictly negative, which implies eventual negativity.  Part~\textup{(iii)}
is immediate from \(K+N\to\infty\).  Thus the sign change cannot occur inside the half-plane
where the positive Dirichlet expansion represents the source.

For part~\textup{(iv)}, use the absolutely convergent identity
\[
 \frac1{\zeta(u)}=\sum_{n\geq1}\frac{\mu(n)}{n^u}
 \qquad(\Re u>1)
\]
and \(\mu(1)=1,\ \mu(2)=-1\).  Finally, if
\(a_r(n)\geq0\) converged coefficientwise to \(b(n)\), then closedness of
\([0,\infty)\) would imply \(b(1)\geq0\), contradicting \(b(1)=-1/2\).
\end{proof}

The theorem is a no-regular-positive-limit principle.  The sign needed for M\"obius inversion
does not appear by cancellation among the positive coefficients: it appears only after the
frequency scale moves, the positive series ceases to represent the function, and meromorphic
continuation is taken.  This is why the reciprocal divisor is a boundary object rather than an
ordinary positive limit.

The selector itself is not specific to zeta.  It is the boundary law of every dilation quotient.

\begin{theorem}[Universal dilation-cocycle law]\label{thm:universal-cocycle}
Let \(F\) be continuous at \(0\), and define
\[
 Q_t^F(s)=\frac{F(s)}{F(ts)}.
 \tag{7.12}
\]
Where the quotients are defined meromorphically,
\[
 Q_{tv}^F(s)=Q_t^F(s)Q_v^F(ts).
 \tag{7.13}
\]
If \(\varepsilon_n\to0\), \(t_n\varepsilon_n\to\lambda\), \(F\) is continuous at
\(\lambda u\), and \(F(\lambda u)\ne0\), then
\[
 Q_{t_n}^F(\varepsilon_nu)
 \longrightarrow\frac{F(0)}{F(\lambda u)}.
 \tag{7.14}
\]
If \(F(0)\ne0\) and \(F\) is meromorphic at \(\rho\), then
\[
 \ord_\rho\frac{F(0)}{F}=-\ord_\rho F.
 \label{eq:order-polarity}
 \tag{7.15}
\]
\end{theorem}

\begin{proof}
The cocycle identity is cancellation of \(F(ts)\).  The limit follows from continuity of the
numerator and denominator.  The order identity is the quotient law for local meromorphic order,
since the nonzero constant \(F(0)\) has order zero.
\end{proof}

Equations \eqref{eq:prefix-transition} and \eqref{eq:order-polarity} explain ``divisor transition''
literally: the same positive finite-depth family tends to the zeta divisor at fixed scale and to
its negative at critical scale.  The product of the two limiting carriers is the endpoint
constant,
\[
 \zeta(u)\frac{\zeta(0)}{\zeta(u)}=\zeta(0),
 \tag{7.16}
\]
away from the divisor.  Divisor order is conserved with opposite polarity.

\section{Reflection polarization and observer loss}\label{sec:polarization}

The arithmetic scale law produces reciprocal zeta, not reciprocal \(\Xi\).  Completion is a
separate operation that removes the trivial divisor and imposes the reflection
\(s\mapsto1-s\).  This distinction prevents the Gamma factor from being smuggled into an
arithmetic source.

\begin{theorem}[Unique affine reflection polarization]\label{thm:affine-polarization}
Let \(F\) be a nonconstant meromorphic function satisfying
\[
 F(1-s)=F(s),\qquad F(0)\ne0.
 \tag{8.1}
\]
Then:
\begin{enumerate}[label=\textup{(\roman*)}]
\item the endpoint charts
\[
 \frac{F(\varepsilon u)}{F(u)}
 \quad\text{and}\quad
 \frac{F(1-\varepsilon u)}{F(u)}
 \tag{8.2}
\]
coincide for every \(\varepsilon,u\), and tend to \(F(0)/F(u)\) as
\(\varepsilon\to0\), away from its poles;
\item an affine function \(w(s)=as+b\) satisfies
\[
 w(1-s)=-w(s)\quad\text{for every }s
 \label{eq:affine-odd}
 \tag{8.3}
\]
if and only if \(b=-a/2\), equivalently \(w(s)=a(s-\tfrac12)\);
\item the polarized boundary current
\begin{equation}
 \mathcal J_F(s)=
 -\left(s-\frac12\right)\frac{F'(s)}{F(s)}
 \label{eq:polarized-current}
\end{equation}
is reflection invariant;
\item if \(\rho\) has meromorphic order \(m\), then
\begin{equation}
 \operatorname*{Res}_{s=\rho}\mathcal J_F
 =-m\left(\rho-\frac12\right).
 \label{eq:polarized-residue}
\end{equation}
\end{enumerate}
\end{theorem}

\begin{proof}
Part~\textup{(i)} follows by substituting \(1-\varepsilon u\) into the reflection identity and
then using continuity at zero.  For part~\textup{(ii)}, setting \(s=\tfrac12\) in
\eqref{eq:affine-odd} gives \(a/2+b=0\); the converse is immediate.

Logarithmic differentiation of \(F(1-s)=F(s)\) gives
\[
 \frac{F'(1-s)}{F(1-s)}=-\frac{F'(s)}{F(s)}.
 \tag{8.6}
\]
Both factors in \eqref{eq:polarized-current} change sign under reflection, proving
part~\textup{(iii)}.  Finally, \(F'/F\) has residue \(m\) at a point of order \(m\);
multiplication by the holomorphic affine weight proves part~\textup{(iv)}.
\end{proof}

Apply the theorem to \(F=\Xi\).  The current does not prove that its residues have zero real
part, but it identifies the unique unsquared coordinate in which that statement can be made.
The familiar observables are projections of the same residue vector.

\begin{corollary}[Projection law]\label{cor:projection-law}
At a zero \(\rho\) of multiplicity \(m\), put
\[
 v_\rho=m(\rho-\tfrac12).
 \tag{8.7}
\]
Then:
\begin{enumerate}[label=\textup{(\roman*)}]
\item logarithmic order retains \(m\) and forgets \(v_\rho/m\);
\item the squared centered coordinate retains
\((\rho-\tfrac12)^2\) but identifies \(\rho\) with \(1-\rho\);
\item the normal rate \(2\Re\rho-1\) is \(2\Re(v_\rho/m)\);
\item a simple off-line functional-equation quartet has argument-principle mass
\(4(2\pi)=8\pi\), while the vector sum of its four centered residues is zero.
\end{enumerate}
No reflection-even scalar observer is injective on \(\C\).
\end{corollary}

\begin{proof}
The first three assertions are the displayed definitions.  Reflection sends
\(\rho-\tfrac12\) to its negative, so its square is unchanged.  A simple quartet contains four
zeros, each contributing \(2\pi\) to the argument-principle mass.  Functional-equation and
Schwarz reflection produce centered vectors
\[
 z,\quad-\overline z,\quad\overline z,\quad-z,
 \tag{8.8}
\]
whose sum is zero.  Finally, if \(O(1-s)=O(s)\), then
\(O(\tfrac12+1)=O(\tfrac12-1)\) at two distinct points, so \(O\) is not injective.
\end{proof}

This gives a precise interpretation of the principal no-go results.  Fixed-scale convergence
keeps the bulk and omits the reciprocal boundary.  Zero renormalized depth collapses the
boundary to the identity.  Right-escaping depth collapses it to a constant.  Finite
differential algebra sees maximally free normalized channels after the common divisor has been
factored out.  Ward or argument-principle mass sees multiplicity but cancels affine orientation
over a quartet.  Squared Julia or Hilbert norms are reflection-even.  None of these failures
contradicts divisor recovery; each is the loss of one coordinate of the transition datum.

\section{A quantitative completed orientation budget}\label{sec:orientation-budget}

The transition law is structural.  We now state the unconditional numerical content supplied by
the completed prime--Gamma source.  Let \(\mathcal D_\Xi\) be the multiset of nontrivial zeros,
counted with analytic multiplicity.  For \(\rho=\beta+i\gamma\), set
\[
 \rho^\dagger=1-\overline\rho,\qquad
 c_\rho=\Re\frac1{\rho(1-\rho)},\qquad
 \lambda_\rho=-\frac4{(2\rho-1)^2}.
 \tag{9.1}
\]
Define
\[
 C=\sum_{\rho\in\mathcal D_\Xi}c_\rho,\qquad
 C_\perp=\sum_{\substack{\rho\in\mathcal D_\Xi\\\Re\rho\ne1/2}}c_\rho,
 \qquad
 \mathcal E_\perp=\sum_{\rho\in\mathcal D_\Xi}(\Im\lambda_\rho)^2.
 \tag{9.2}
\]
The sums converge absolutely.  On-line terms in \(\mathcal E_\perp\) vanish.

\begin{lemma}[Positive endpoint charge]\label{lem:endpoint-charge}
For every \(\rho=\beta+i\gamma\in\mathcal D_\Xi\),
\[
 c_\rho=
 \frac{\gamma^2+\beta(1-\beta)}
 {(\beta^2+\gamma^2)((1-\beta)^2+\gamma^2)}>0.
 \label{eq:endpoint-charge-formula}
 \tag{9.3}
\]
Moreover
\[
 C=2+\gamma_{\!E}-\log(4\pi)
 <\frac{93}{2000}.
 \tag{9.4}
\]
\end{lemma}

\begin{proof}
Multiplication by the conjugate denominator gives \eqref{eq:endpoint-charge-formula}; positivity uses
\(0<\beta<1\).  The centered Hadamard logarithmic-derivative identity at the completed endpoint
\(s=0\) gives \(C=\sum_\rho c_\rho\).  Expanding
\[
 \Xi(s)=\frac12s(s-1)\pi^{-s/2}\Gamma(s/2)\zeta(s)
 \tag{9.5}
\]
at zero yields the exact value in \eqref{eq:9.4}.  The rational upper bound follows from the
standard explicit inequalities
\(\gamma_{\!E}<2887/5000\) and
\(\log(4\pi)>25309/10000\), which give a value strictly below \(93/2000\).
\end{proof}

\begin{theorem}[Four-orbit tax and cubic orientation law]
\label{thm:orientation-budget}
One has
\[
 0\leq C_\perp\leq C,\qquad
 C_\perp=0\Longleftrightarrow\mathcal E_\perp=0
 \Longleftrightarrow\mathrm{RH}.
 \label{eq:orientation-zero}
 \tag{9.6}
\]
For every off-line zero \(\rho=\beta+i\gamma\),
\[
 4c_\rho\leq C_\perp,\qquad
 4\leq C_\perp(\gamma^2+1),\qquad
 \gamma^2>\frac{7907}{93}.
 \label{eq:height-gap}
 \tag{9.7}
\]
Globally,
\begin{equation}
 \mathcal E_\perp
 \leq4\left(\frac{C_\perp}{4-C}\right)^3
 \leq4\left(\frac{C}{4-C}\right)^3
 <\frac1{150000}.
 \label{eq:orientation-cubic}
\end{equation}
\end{theorem}

\begin{proof}
Positivity gives \(0\leq C_\perp\leq C\).  Both \(C_\perp\) and
\(\mathcal E_\perp\) are sums of nonnegative terms, and their individual terms vanish exactly
when \(\beta=\tfrac12\), proving \eqref{eq:orientation-zero}.

If \(\rho\) is off line, the four points
\[
 \rho,\quad1-\overline\rho,\quad\overline\rho,\quad1-\rho
 \tag{9.9}
\]
are distinct and have the same positive charge.  Hence \(4c_\rho\leq C_\perp\).  Formula
\eqref{eq:endpoint-charge-formula} and \(0<\beta<1\) give
\[
 c_\rho\geq\frac1{\gamma^2+1}.
 \tag{9.10}
\]
Thus \(4\leq C_\perp(\gamma^2+1)\leq C(\gamma^2+1)\).  Using
\(C<93/2000\) gives
\[
 \gamma^2>\frac4C-1>\frac{8000}{93}-1
 =\frac{7907}{93}.
 \tag{9.11}
\]

For the cubic estimate, put
\[
 x=(\beta-\tfrac12)^2,\qquad v=\gamma^2.
 \]
Then \(0\leq x\leq1/4\), and direct calculation gives
\[
 (\Im\lambda_\rho)^2
 =\frac{4xv}{(x+v)^4}\leq\frac1{v^3}.
 \tag{9.12}
\]
The inequality \(4\leq C(v+1)\) is equivalent to
\[
 \frac1v\leq\frac4{4-C}\frac1{v+1}.
 \tag{9.13}
\]
Since \((v+1)^{-1}\leq c_\rho\), cubing gives
\[
 (\Im\lambda_\rho)^2
 \leq\left(\frac4{4-C}\right)^3c_\rho^3.
 \label{eq:atom-cubic}
 \tag{9.14}
\]
For an off-line term, the quartet bound gives
\(c_\rho\leq C_\perp/4\), hence
\[
 c_\rho^3\leq(C_\perp/4)^2c_\rho.
 \tag{9.15}
\]
Summing \eqref{eq:atom-cubic} and using
\(\sum_{\Re\rho\ne1/2}c_\rho=C_\perp\) proves the first inequality in
\eqref{eq:orientation-cubic}.  Monotonicity in \(C_\perp\) proves the second.  Finally,
substitution of \(C<93/2000\) and exact rational cross-multiplication gives
\[
 4\left(\frac{93/2000}{4-93/2000}\right)^3
 <\frac1{150000},
 \]
which proves the last inequality.
\end{proof}

\begin{corollary}[Exact source gap]\label{cor:source-gap}
Let
\[
 M_2=\sum_{\rho\in\mathcal D_\Xi}(\Re\lambda_\rho)^2
\]
and let \(\mathcal Q_\Gamma\) be the second centered moment of the completed prime--Gamma
source.  Then
\[
 M_2-\mathcal Q_\Gamma=\mathcal E_\perp,\qquad
 0\leq M_2-\mathcal Q_\Gamma<\frac1{150000}.
 \tag{9.17}
\]
Equality \(M_2=\mathcal Q_\Gamma\) is equivalent to RH.
\end{corollary}

\begin{proof}
The completed logarithmic derivative has the locally convergent expansion
\[
 -\frac1z\frac{d}{dz}\log\Xi(\tfrac12+iz)
 =\sum_\rho\frac{\lambda_\rho}{1-z^2\lambda_\rho}.
 \tag{9.18}
\]
Its second source moment is the real part of
\(\sum_\rho\lambda_\rho^2\).  Since
\[
 \Re(\lambda_\rho^2)
 =(\Re\lambda_\rho)^2-(\Im\lambda_\rho)^2,
 \]
absolute summability gives
\(M_2-\mathcal Q_\Gamma=\mathcal E_\perp\).  Apply
Theorem~\ref{thm:orientation-budget}; equality is equivalent to
\(\mathcal E_\perp=0\).
\end{proof}

\section{Closed consequences for RH routes}\label{sec:closed-scope}

Theorems~\ref{thm:complete-frame}, \ref{thm:scale-selector},
\ref{thm:positive-mobius}, \ref{thm:affine-polarization}, and
\ref{thm:orientation-budget} give the following closed structural conclusions.

\begin{enumerate}[label=\textup{(\arabic*)}]
\item The divisor is not encoded by overlap of coefficient supports or by a finite differential
relation.  It is longitudinal data shared by a maximally transverse positive resolution.
\item A divisor survives a vanishing-coordinate limit precisely in the finite nonzero
renormalized-depth class covered by Theorem~\ref{thm:scale-selector}.  Zero depth gives the
identity; right escape gives a constant.
\item Signed M\"obius inversion cannot arise as regular coefficientwise convergence of the
positive sources.  Divisor retention forces moving-frequency scale escape and analytic
continuation.
\item The exceptional boundary reverses divisor order.  Thus positivity in the bulk and
negative divisor polarity at the boundary are compatible because they belong to different
charts.
\item Hypertranscendence cannot itself exclude zeros: it is transverse to the common zeta
factor.  Maximal differential rank and complete common-divisor replication coexist.
\item Ward quartet counts retain multiplicity but cancel the affine residue vector; squared
Julia coordinates and reflection-even norms identify the two horizontal sides.  The centered
affine current is the unique affine reflection-compatible carrier of that side.
\item Every hypothetical off-line orbit pays four positive endpoint charges, obeys the height
bound \eqref{eq:height-gap}, and contributes to a total orientation energy smaller than
\(1/150000\).
\end{enumerate}

No open hypothesis is hidden in these conclusions.  They do not assert that the orientation
energy vanishes and therefore do not assert RH.  Their unconditional content is the exact
source-to-boundary mechanism, its necessity, its information-loss classification, and its
quantitative defect budget.  Any method that keeps only a fixed-scale jet, a finite differential
model, a zero count, or a squared coordinate has already discarded one of the coordinates in
that mechanism.

\section{Conclusion}

One explicit arithmetic object realizes the full law.  The complete largest-exponent resolution
is positive and pairwise disjoint; its normalized channels have maximal differential rank over
the zeta field; every channel carries the same divisor; and its exact finite prefixes form a
dilation cocycle.  At native real points those prefixes are positive.  At finite nonzero
renormalized depth they leave the positive Dirichlet domain, become eventually negative, and
converge to a mixed-sign M\"obius inverse.  At that boundary, zeros become poles and every local
divisor order reverses.

Archimedean completion does not create this transition.  It polarizes the divisor already
exposed by arithmetic depth.  Reflection selects the centered affine coordinate uniquely, and
the corresponding logarithmic current has the full vector residue
\(-m(\rho-\tfrac12)\).  Counts, squares, norms, compact limits, and finite differential fields
are then understood as projections rather than failed reconstructions.

The completed prime--Gamma source adds strict quantitative content.  An off-line quartet has a
four-charge cost, every off-line zero obeys
\(\gamma^2>7907/93\), and the total transverse orientation energy is below \(1/150000\).
RH is exactly the zero-energy case of that already source-bounded polarized defect.  This paper
proves the sign-transmutation law, the divisor-polarity law, the unique affine polarization, and
the quantitative budget unconditionally; it neither assumes nor announces RH.

\begin{thebibliography}{99}

\bibitem{GrosswaldSchnitzer1978}
E.~Grosswald and F.~J.~Schnitzer,
\newblock A class of modified \(\zeta\) and \(L\)-functions,
\newblock \emph{Pacific J. Math.} \textbf{74} (1978), 357--364.
\newblock \href{https://doi.org/10.2140/pjm.1978.74.357}
{doi:10.2140/pjm.1978.74.357}.

\bibitem{Kawalec2026}
A.~Kawalec,
\newblock On the series expansion of \(k\)-free Dirichlet series and its analytical continuation,
\newblock \emph{arXiv:2603.22775} (2026).
\newblock \href{https://arxiv.org/abs/2603.22775}{arXiv:2603.22775}.

\bibitem{Kolchin1973}
E.~R.~Kolchin,
\newblock \emph{Differential Algebra and Algebraic Groups},
\newblock Academic Press, New York, 1973.

\bibitem{Nagoshi2016}
H.~Nagoshi,
\newblock The hypertranscendence of \(L\)-functions for
\(\mathrm{GL}_m(\mathbb A_{\mathbb Q})\),
\newblock \emph{Bull. Aust. Math. Soc.} \textbf{93} (2016), 388--399.
\newblock \href{https://doi.org/10.1017/S000497271500129X}
{doi:10.1017/S000497271500129X}.

\bibitem{Meng2015}
X.~Meng,
\newblock The distribution of \(k\)-free numbers and the derivative of the Riemann
\(\zeta\)-function,
\newblock \emph{Math. Proc. Cambridge Philos. Soc.} \textbf{162} (2017), 293--317.
\newblock \href{https://doi.org/10.1017/S0305004116000554}
{doi:10.1017/S0305004116000554}.

\bibitem{Ostrowski1920}
A.~Ostrowski,
\newblock \"Uber Dirichletsche Reihen und algebraische Differentialgleichungen,
\newblock \emph{Math. Z.} \textbf{8} (1920), 241--298.
\newblock \href{https://doi.org/10.1007/BF01206523}
{doi:10.1007/BF01206523}.

\bibitem{Pong2016}
W.~Y.~Pong,
\newblock Applications of differential algebra to algebraic independence of arithmetic
functions,
\newblock \emph{Acta Arith.} \textbf{172} (2016), 149--173.
\newblock \href{https://doi.org/10.4064/aa8112-12-2015}
{doi:10.4064/aa8112-12-2015}.

\bibitem{Postnikov1949}
A.~G.~Postnikov,
\newblock On the differential independence of Dirichlet series,
\newblock \emph{Dokl. Akad. Nauk SSSR} \textbf{66} (1949), 561--564
\newblock (in Russian).

\bibitem{ShapiroSparer1986}
H.~N.~Shapiro and G.~H.~Sparer,
\newblock On algebraic independence of Dirichlet series,
\newblock \emph{Comm. Pure Appl. Math.} \textbf{39} (1986), 695--745.
\newblock \href{https://doi.org/10.1002/cpa.3160390602}
{doi:10.1002/cpa.3160390602}.

\bibitem{Titchmarsh1986}
E.~C.~Titchmarsh,
\newblock \emph{The Theory of the Riemann Zeta-Function}, 2nd ed.,
revised by D.~R.~Heath-Brown,
\newblock Oxford University Press, Oxford, 1986.

\bibitem{Voronin1973}
S.~M.~Voronin,
\newblock On differential independence of \(\zeta\)-functions,
\newblock \emph{Soviet Math. Dokl.} \textbf{14} (1973), 607--609.
\newblock \href{https://www.mathnet.ru/eng/dan37604}{MathNet publication page}.

\end{thebibliography}

\end{document}
